\documentclass{article}
\usepackage{cmsa-abs}
\begin{document}

\abs Leopardi %% put speaker's surname here
\sp Paul Leopardi %% put speaker's whole name here
\au Paul Leopardi %% List all authors here
\aff Australian National University %% Give speaker's affiliation here
\ams 05Z99 %% List Mathematics Subject Classification here

\title Big data in combinatorics: is it FAIR?

\text 
In 2015, a Lorentz Center workshop drafted the FAIR 
(Findable, Accessible, Interoperable, Reusable) Data Principles for research data.
Commenting on a 2016 meeting at St. Andrews on \emph{Discrete Mathematics and Big Data}, Peter Cameron stated that lists of generated data, for example, lists of semigroups, need to be stored in accessible form.
The EU funded OpenDreamKit project (2015-2019) included a census of mathematical data sets, and infrastructure for the FAIR sharing of such data sets.

This talk will examine some combinatorial data sets from the FAIR perspective, and describe work in progress to share a data set containing classifications of bent functions by their Cayley graphs.

\endabs

\end{document}
